% Options for packages loaded elsewhere
\PassOptionsToPackage{unicode}{hyperref}
\PassOptionsToPackage{hyphens}{url}
\PassOptionsToPackage{dvipsnames,svgnames,x11names}{xcolor}
%
\documentclass[
  10pt,
]{article}
\usepackage{amsmath,amssymb}
\usepackage{iftex}
\ifPDFTeX
  \usepackage[T1]{fontenc}
  \usepackage[utf8]{inputenc}
  \usepackage{textcomp} % provide euro and other symbols
\else % if luatex or xetex
  \usepackage{unicode-math} % this also loads fontspec
  \defaultfontfeatures{Scale=MatchLowercase}
  \defaultfontfeatures[\rmfamily]{Ligatures=TeX,Scale=1}
\fi
\usepackage{lmodern}
\ifPDFTeX\else
  % xetex/luatex font selection
  \setmainfont[]{DejaVu Serif}
  \setsansfont[]{DejaVu Sans}
  \setmonofont[]{DejaVu Sans Mono}
  \setmathfont[]{latinmodern-math.otf}
\fi
% Use upquote if available, for straight quotes in verbatim environments
\IfFileExists{upquote.sty}{\usepackage{upquote}}{}
\IfFileExists{microtype.sty}{% use microtype if available
  \usepackage[]{microtype}
  \UseMicrotypeSet[protrusion]{basicmath} % disable protrusion for tt fonts
}{}
\makeatletter
\@ifundefined{KOMAClassName}{% if non-KOMA class
  \IfFileExists{parskip.sty}{%
    \usepackage{parskip}
  }{% else
    \setlength{\parindent}{0pt}
    \setlength{\parskip}{6pt plus 2pt minus 1pt}}
}{% if KOMA class
  \KOMAoptions{parskip=half}}
\makeatother
\usepackage{xcolor}
\usepackage[a4paper,margin=24mm]{geometry}
\setlength{\emergencystretch}{3em} % prevent overfull lines
\providecommand{\tightlist}{%
  \setlength{\itemsep}{0pt}\setlength{\parskip}{0pt}}
\setcounter{secnumdepth}{-\maxdimen} % remove section numbering
\ifLuaTeX
\usepackage[bidi=basic]{babel}
\else
\usepackage[bidi=default]{babel}
\fi
\babelprovide[main,import]{british}
\ifPDFTeX
\else
\babelfont{rm}[]{DejaVu Serif}
\fi
% get rid of language-specific shorthands (see #6817):
\let\LanguageShortHands\languageshorthands
\def\languageshorthands#1{}
\usepackage{fancyhdr}
\usepackage{microtype}
\usepackage{needspace}
\pagestyle{fancy}
\fancyhf{}
\fancyhead[L]{\small\sffamily PROPOSED RESOLUTION / KE-04}
\fancyhead[R]{\small\sffamily 11 September 2026}
\fancyfoot[L]{\footnotesize Not independently verified}
\fancyfoot[R]{\footnotesize \thepage}
\setlength{\headheight}{15pt}
\setlength{\emergencystretch}{3em}
\allowdisplaybreaks[1]
\makeatletter
\renewcommand{\maketitle}{%
  \thispagestyle{fancy}%
  \begin{flushleft}
  {\Large\bfseries \@title\par}
  \vspace{0.5em}
  {\normalsize\sffamily Proposed resolution for mathematical review\par}
  \end{flushleft}
  \vspace{0.5em}
}
\makeatother
\ifLuaTeX
  \usepackage{selnolig}  % disable illegal ligatures
\fi
\usepackage{bookmark}
\IfFileExists{xurl.sty}{\usepackage{xurl}}{} % add URL line breaks if available
\urlstyle{same}
\hypersetup{
  pdftitle={KE-04: Strict interval occupancy across block Lanczos iterations},
  pdflang={en-GB},
  colorlinks=true,
  linkcolor={Maroon},
  filecolor={Maroon},
  citecolor={Blue},
  urlcolor={blue},
  pdfcreator={LaTeX via pandoc}}

\title{KE-04: Strict interval occupancy across block Lanczos iterations}
\author{}
\date{11 September 2026}

\begin{document}
\maketitle

\textbf{Status of this manuscript:} Proposed resolution, not
independently verified.\\
\textbf{Outcome claimed:} Affirmative resolution claim.\\
\textbf{Prepared:} 11 September 2026.\\
\textbf{Target:}
\href{https://github.com/ajt60gaibb/OpenProblemsInNLA/blob/b4123194697bdf6f8f82518c1dd7d6c40a30c2e0/eigenvalues-and-inverse-problems/KE-04/README.md}{Repository
entry KE-04}, snapshot \texttt{b412319}.

This is a proposed argument generated in a ChatGPT conversation and
prepared for mathematical review. No independent referee report or
formal proof certificate accompanies it. Supporting algebraic and
numerical diagnostics are not a substitute for proof review. No
publication or priority claim is made.

\subsection{Theorem KE-04: proposed strict interval
occupancy}\label{theorem-ke-04-proposed-strict-interval-occupancy}

Let \(A\in\mathbb R^{n\times n}\) be symmetric and let
\(V\in\mathbb R^{n\times p}\) have full column rank. Define \[
\mathcal K_\ell=\operatorname{range}[V,AV,\ldots,A^{\ell-1}V].
\] Assume \(\dim\mathcal K_s=sp\) and therefore full block dimension at
all preceding iterations. Let \(P_\ell\) be the orthogonal projector
onto \(\mathcal K_\ell\), and let
\(H_\ell=P_\ell A|_{\mathcal K_\ell}\). Write its eigenvalues, with
multiplicity, as \[
\theta_1^{(\ell)}\le\cdots\le\theta_{\ell p}^{(\ell)}.
\] These are the Ritz values and are independent of the chosen
orthonormal basis.

\textbf{Claim.} For every \(1\le k<j\le s\) and every
\(1\le i\le(k-1)p\), the open interval \[
(\theta_i^{(k)},\theta_{i+p}^{(k)})
\] contains an eigenvalue of \(H_j\).

\subsubsection{1. Assume a gap at the later
iteration}\label{assume-a-gap-at-the-later-iteration}

There are no admissible indices when \(k=1\), so suppose \(k\ge2\). Put
\[
a=\theta_i^{(k)},\qquad b=\theta_{i+p}^{(k)},\qquad
q(t)=(t-a)(t-b).
\] Suppose that \(H_j\) has no eigenvalue in \((a,b)\). Since
\(a\le b\), \[
q(H_j)\succeq0.
\] Let \(E\subseteq\mathcal K_k\) be the span of an orthonormal set of
eigenvectors of \(H_k\) corresponding to indices \(i,\ldots,i+p\). Then
\[
\dim E=p+1,\qquad \langle x,q(H_k)x\rangle\le0\quad(x\in E).
\] The codimension of \(\mathcal K_{k-1}\) in \(\mathcal K_k\) is \(p\),
so \(E\cap\mathcal K_{k-1}\) contains a nonzero vector \(x\).

\subsubsection{2. The quadratic forms agree on that
vector}\label{the-quadratic-forms-agree-on-that-vector}

Since \(x\in\mathcal K_{k-1}\), we have \(Ax\in\mathcal K_k\) and \[
H_kx=H_jx=Ax.
\] Self-adjointness of the compressions gives \[
\begin{aligned}
\langle x,q(H_j)x\rangle
&=\|Ax\|^2-(a+b)\langle x,Ax\rangle+ab\|x\|^2\\
&=\langle x,q(H_k)x\rangle\le0.
\end{aligned}
\] Combined with \(q(H_j)\succeq0\), this implies \(q(H_j)x=0\).

Moreover \(A^2x\in\mathcal K_{k+1}\subseteq\mathcal K_j\), so
\(H_j^2x=A^2x\). Hence \[
q(A)x=0.
\]

\subsubsection{3. Full block rank excludes this
annihilation}\label{full-block-rank-excludes-this-annihilation}

Write \[
x=\sum_{\ell=0}^{k-2}A^\ell Vc_\ell.
\] Let \(r\) be the largest index with \(c_r\ne0\); such an index exists
because \(x\ne0\). The polynomial \(q\) is monic of degree two.
Therefore the coefficient of \(A^{r+2}V\) in the resulting block-power
expansion of \(q(A)x\) is exactly \(c_r\ne0\).

Since \(k+1\le j\le s\), the block Krylov matrix \[
[V,AV,\ldots,A^kV]
\] has full column rank. As \(r+2\le k\), the nonzero coefficient just
found implies \(q(A)x\ne0\), a contradiction.

This establishes the claim for all admissible indices and iterations. If
\(a=b\), the same argument applies with \(q(t)=(t-a)^2\) and rules out
that case; thus the proof does not presuppose strictly separated
endpoints or simple Ritz values. \(\square\)

\subsection{Scope and review notes}\label{scope-and-review-notes}

The argument uses real symmetric A, a full-column-rank starting block,
exact arithmetic, and full block dimension through iteration j. It
includes eigenvalue multiplicities and the possibility of coincident
proposed interval endpoints, which it rules out in the relevant range.
It does not assert a finite-precision analogue.

The key review point is the equality of the two quadratic forms on a
vector in \(\mathcal K_{k-1}\), followed by conversion to a degree-two
annihilating polynomial. Full rank is required through
\(\mathcal K_{k+1}\), which is guaranteed by \(k+1\le j\le s\).

The packaging pass checked the repository statement and contribution
rules on 11 September 2026. It did not conduct a new exhaustive
literature search or establish novelty.

\subsection{Source references}\label{source-references}

Alex Townsend, \emph{Open Problems in Numerical Linear Algebra} (2026),
\href{https://github.com/ajt60gaibb/OpenProblemsInNLA/blob/b4123194697bdf6f8f82518c1dd7d6c40a30c2e0/eigenvalues-and-inverse-problems/KE-04/README.md}{entry
KE-04}, repository snapshot \texttt{b412319}, accessed 11 September
2026.

D. Šimonová and P. Tichý, \emph{On finite precision block Lanczos
computations}, \href{https://arxiv.org/abs/2507.16484v1}{arXiv
manuscript, version 1}, 22 July 2025, §2.1, unnumbered conjecture. The
catalog identifies this as the source of its exact-arithmetic interval
statement.

\end{document}
